\documentclass[11pt]{article}
\usepackage[margin=1in]{geometry}
\usepackage{amsmath,amssymb,amsthm,hyperref,enumitem}
\hypersetup{colorlinks=true,linkcolor=blue,urlcolor=blue}
\newtheorem{theorem}{Theorem}
\newtheorem{proposition}{Proposition}
\newcommand{\Corr}{\operatorname{Corr}}
\title{The Euler-transform candidate need not optimize\\the generating-function correlation}
\author{Principia Math In Progress}
\date{22 September 2026}
\begin{document}
\maketitle
\noindent\textbf{MathDB 356222. Full disproof of Conjecture 5.93.}

\section*{Task 1: the intended optimization problem}
For an invertible lower-triangular kernel \(D\), Schmidt [1,
Notation 5.85 and Question 5.86, pp.~79--80] considers
\[
 \Corr(C,D)=\sum_{n\geq1}\frac1n
 \frac{\sum_{k=1}^n|c_k D^{-1}(n,k)|}
 {\left(\sum_{k=1}^nc_k^2\right)^{1/2}
  \left(\sum_{k=1}^nD^{-1}(n,k)^2\right)^{1/2}},
 \qquad c_k=[q^k]C(q).                                  \tag{1}
\]
The stated optimization ranges over rational formal power series with
\(C(0)\ne0\), or normalized integer series with \(C(0)=1\), and
requires a finite score. Conjecture 5.93 [1, p.~87] proposes
the Euler-transform expression
\[
 C_D(q)=\exp\left(-\sum_{n\geq1}
                    \frac{q^n}{n}\sum_{k=1}^n kD(n,k)\right)       \tag{2}
\]
as an optimal selection. We disprove its actual optimality:
an admissible kernel gives a well-defined candidate with score
strictly between the scores of two admissible integer polynomials.
Thus the counterexample is independent of the additional product
formula printed alongside (2).

\begin{theorem}
Let \(D(n,k)=\delta_{nk}\), so \(D^{-1}=D\), and write
\[
 E(q)=\exp\left(-\frac{q}{1-q}\right),\qquad
 P_N(q)=1+\sum_{k=1}^N2^{k-1}q^k .
\]
Every \(P_N\) has a finite score, and
\[
 \Corr(P_N,D)>\frac{\sqrt3}{2}H_N,\qquad
 H_N=\sum_{n=1}^N\frac1n.                                \tag{3}
\]
Consequently, finite scores are unbounded even on normalized
integer polynomials, so no finite global maximum exists.
The proposed candidate (2) is \(E\), belongs to
\(\mathbb Q[[q]]\), and satisfies
\[
 1=\Corr(1+q,D)<\Corr(E,D)<2
      <\frac{49\sqrt3}{40}<\Corr(P_6,D).                 \tag{4}
\]
In particular, \(E\) is neither a maximum nor a minimum.
\end{theorem}

\begin{proof}
The identity kernel is lower triangular, invertible, integral, and
has diagonal entries one. The inner sum in (2) is \(n\), giving
\(C_D=E\). Its constant and linear coefficients are \(1,-1\);
its coefficients are rational by the formal exponential expansion.
For any real coefficients with \(c_1\ne0\), (1) becomes
\[
 F(C):=\Corr(C,D)
   =\sum_{n\geq1}\frac{|c_n|}
      {n\left(\sum_{k=1}^n c_k^2\right)^{1/2}}.           \tag{5}
\]
All denominators used below are positive, including those in the
zero tail of a polynomial. The first summand is one.

For \(P_N\), only the first \(N\) summands can be nonzero and
\[
 \sum_{k=1}^n c_k^2=\frac{4^n-1}{3},\qquad
 \frac{c_n}{\left(\sum_{k=1}^nc_k^2\right)^{1/2}}
 =\frac{\sqrt3}{2\sqrt{1-4^{-n}}}>\frac{\sqrt3}{2}
 \quad(1\leq n\leq N).
\]
This proves (3). Since the harmonic sums are unbounded, it also
proves the nonexistence of a finite maximum, with every individual
competitor satisfying the source's finiteness requirement.

It remains to prove the finite upper bound for \(E\), without
assuming any asymptotic formula for its coefficients.
For \(|z|<1\),
\[
 \Re\frac{z}{1-z}
    =-\frac12+\frac{1-|z|^2}{2|1-z|^2}>-\frac12,
 \qquad |E(z)|^2<e.                                    \tag{6}
\]
Write \(E(z)=\sum_{n\geq0}e_nz^n\). For \(0<t<1\), uniform
convergence of the power series on \(|z|=t\), followed by
orthogonality of exponentials, gives
\[
 \sum_{n\geq0}e_n^2t^{2n}
 =\frac1{2\pi}\int_0^{2\pi}|E(te^{i\theta})|^2\,d\theta
 \leq e.
\]
Letting \(t\uparrow1\), monotone convergence and \(e_0=1\) show
\(\sum_{n\geq1}e_n^2\leq e-1<2\).
Here \(e<3\) follows directly from its factorial series.
Also
\[
 \sum_{n\geq1}\frac1{n^2}
 <1+\sum_{n\geq2}\frac1{n(n-1)}=2.
\]
Since \(e_1=-1\), every coefficient norm in (5) is at least one.
Cauchy--Schwarz therefore proves absolute convergence and
\[
 F(E)\leq\sum_{n\geq1}\frac{|e_n|}{n}
 \leq\left(\sum_{n\geq1}e_n^2\right)^{1/2}
     \left(\sum_{n\geq1}n^{-2}\right)^{1/2}<2.            \tag{7}
\]
The expansion \(E(q)=1-q-\tfrac12q^2+O(q^3)\) implies
\(F(E)\geq1+1/(2\sqrt5)>1\).
Finally, \(H_6=49/20\), and
\((49\sqrt3)^2=7203>6400=80^2\). Equations (3) and (7)
give all the strict inequalities in (4).
\end{proof}

\section*{A sharp finite-degree statement}
\begin{proposition}
For \(N\geq1\), the supremum of (5) over normalized integer
polynomials of degree at most \(N\), with \(c_1\ne0\), is \(H_N\).
For \(N\geq2\) the supremum is not attained. The minimum is one,
attained precisely when all \(c_n=0\) for \(n\geq2\).
\end{proposition}
\begin{proof}
Each summand in (5) is at most \(1/n\). For \(n=2\) it is strictly
less than \(1/2\) because \(c_1\ne0\), even if \(c_2=0\).
For integers \(R\geq2\), the normalized integer polynomial
\(1+\sum_{k=1}^N R^{k-1}q^k\) has \(n\)-th coefficient ratio
\[
 \frac{R^{n-1}}{\sqrt{1+R^2+\cdots+R^{2n-2}}}
   =\sqrt{\frac{1-R^{-2}}{1-R^{-2n}}}\longrightarrow1 .
\]
The finite sum tends to \(H_N\). The minimum characterization
follows from the first summand being one and all others nonnegative.
\end{proof}

\section*{Scope and attribution}
The source explicitly allows all invertible lower-triangular kernels
and finite-score rational OGFs; the identity kernel and all
competitors above satisfy those conditions.
No divergent-score series is used as a competitor.
The printed infinite-product expression has a separate formal
convergence issue; neither that issue nor a change in normalization
is the basis of this disproof. The power series (2) is actually
outperformed by a normalized integer polynomial.
The specialized Lambert-series optimization problem is not resolved
by this identity-kernel example. The common obstruction to general
maximizer existence also bears on neighboring proposals, but only
the selected Conjecture 5.93 receives a campaign credit.

\section*{Tasks 2--5: proposed mathematical follow-ups}
\begin{enumerate}[label=\textbf{Task \arabic*:},start=2,leftmargin=*]
\item Characterize invertible lower-triangular kernels \(D\) for which
the finite values of (1) are uniformly bounded over normalized
integer polynomials with nonzero linear coefficient.
\item For the Lambert kernel \(D(n,k)=\mathbf1_{k\mid n}\), determine
whether a finite maximum exists over a precisely specified class
of rational OGFs, and whether \((q;q)_\infty\) attains it.
\item With fixed degree \(N\), \(c_1=1\), and coefficient bounds
\(|c_k|\leq B\), determine the exact maximizers of the identity-kernel
score and their behavior as \(N\) or \(B\) grows.
\item Replace \(1/n\) in (1) by a positive summable weight sequence,
and find coefficient-growth constraints giving both existence
of a maximizer and a useful Euler-transform characterization.
\end{enumerate}

\begin{thebibliography}{9}
\bibitem{schmidt}
M.~D. Schmidt, \emph{Factorization theorems and canonical representations
for generating functions of special sums}, Georgia Institute of
Technology doctoral dissertation, August 2022, Chapter~5,
pp.~78--89; \href{https://arxiv.org/abs/2209.12287}{arXiv:2209.12287v1}.
Official final dissertation:
\href{https://repository.gatech.edu/items/f6054390-aaa2-465a-85c2-403619d773ab}
{Georgia Tech repository}.
\end{thebibliography}
\end{document}
